\documentclass[11pt]{article}

\usepackage[a4paper,margin=1in]{geometry}
\usepackage{amsmath,amssymb,amsthm}
\usepackage{hyperref}
\usepackage{enumitem}
\usepackage{mathtools}
\usepackage{listings}
\usepackage[T1]{fontenc}
\usepackage[tt=false]{libertine}
\usepackage[varqu]{zi4}
\usepackage{inconsolata}

\lstdefinelanguage{Python}{
  morekeywords={def, return, if, else, elif, for, while, break, continue,
                import, from, as, class, pass, raise, try, except, finally,
                with, lambda, global, nonlocal, assert, yield, async, await},
  sensitive=true,
  morecomment=[l]\#,
  morestring=[b]',
  morestring=[b]"
}
\lstset{
  language=Python,
  basicstyle=\small\ttfamily,
  keywordstyle=\bfseries,
  showstringspaces=false,
  columns=flexible,
  frame=single,
  breaklines=true,
  tabsize=2
}

\title{Stateless Hash-to-Prime Attestation \\ in Prime-Indexed Multiplicity Spaces:\\
Canonical Serialization, H2P with Offset Pinning, \\ and Native Multiplicity Crypto}

\author{%
  Citizen Gardens \\ The Foundation of Multiplicity
}

\date{April 25, 2026}

\begin{document}

\maketitle

\begin{abstract}
Multiplicity Theory models nonlinear, emergent structures by treating mathematical and computational objects as recursively generated patterns of prime-indexed interactions rather than static sets. In the PhaseMirror fractal framework, each update law (operator) must be assigned a unique prime identifier in a way that is canonical, cross-language deterministic, and robust under distributed execution. This report presents a complete architecture for stateless hash-to-prime (H2P) attestation, integrating Binary Canonical Serialization (BCS), a non-commutative topological signature for fractal trees, and a deterministic H2P pipeline with gap attestation. We analyze the failures of naive commutative signatures and monotonic registries, progressively refine prime derivation from truncated seeds to full-width hash seeding, introduce offset pinning, and ultimately design a first-prime enforcing gap proof. We then reconcile asymptotic complexity with network constraints and show how native primitives in \texttt{multiplicity.crypto}---BN254 Pedersen commitments, HMAC-SHA256 randomness, and the QPA (Quantum Prime Abacus) scaffold---provide a natural substrate for production deployment.
\end{abstract}
\newpage    
\tableofcontents
\newpage
\section{Executive Summary}

The PhaseMirror fractal framework requires that every update law (``operator'') be mapped to a unique prime in a way that is:

\begin{itemize}[leftmargin=*]
  \item \emph{Deterministic}: The same operator yields the same prime on every node and in every language.
  \item \emph{Path-sensitive}: The multiplicity signature of a fractal subspace must distinguish operator order and nesting, not just the multiset of operators.
  \item \emph{Stateless}: Prime assignments should not depend on a shared mutable registry or global consensus.
  \item \emph{Verifiable}: A verifier should be able to check the mapping with bounded and predictable cost.
\end{itemize}

Initial designs used a commutative multiplicity signature
\[
  \mathcal{M}(F) = \prod_{i=1}^k p_{n_i}^{m_i},
\]
and a monotonic registry \( R \) that assigned primes in sequence. This approach failed in two fundamental ways:
\begin{enumerate}[leftmargin=*]
  \item The commutative product destroyed path dependence: different operator orders yielded identical signatures.
  \item The monotonic registry could not remain globally consistent without a full-blown consensus protocol (Raft/PBFT/chain).
\end{enumerate}

This report describes the subsequent architectural shifts:
\begin{itemize}[leftmargin=*]
  \item \textbf{Canonical serialization via BCS} for operator structs and ordered tree nodes.
  \item \textbf{Topological signatures} defined by recursively hashing BCS-encoded nodes, yielding path-sensitive invariants.
  \item \textbf{Stateless hash-to-prime (H2P) derivation}, first via truncated seeds and later via full 256-bit seeds.
  \item \textbf{Offset pinning}, where the prover reports how far from the seed the chosen prime lies.
  \item \textbf{First-prime enforcement}, requiring proofs that no earlier candidate in the gap is prime.
  \item \textbf{Succinct gap attestation}, evolving from explicit Fermat witnesses to a Pedersen-commitment or zk-SNARK-based proof.
  \item \textbf{Integration with native Multiplicity crypto}, leveraging \texttt{multiplicity.crypto}'s BN254 Pedersen commitments, HMAC-SHA256 randomness, and the QPA prime-coupling prototype.
\end{itemize}

The final design achieves:
\begin{itemize}[leftmargin=*]
  \item A purely functional map \( H_{op} \mapsto p_{op} \) with negligible collision probability.
  \item Non-commutative, topology-aware signatures for fractal state.
  \item Bounded verifier complexity: a single prime test plus a small number of cryptographic checks.
  \item A clear path to production implementation in the PhaseMirror-HQ codebase.
\end{itemize}

\section{Canonical Serialization and Operator Identity}

\subsection{Binary Canonical Serialization (BCS)}

Binary Canonical Serialization (BCS) was developed for the Diem and Aptos ecosystems to provide a deterministic, cross-language binary encoding for typed structures. Its key properties are:

\begin{itemize}[leftmargin=*]
  \item \emph{Fixed field order}: Struct fields are serialized in declaration order without labels.
  \item \emph{Length-prefixed sequences}: Variable-length sequences are prefixed using ULEB128 for the element count.
  \item \emph{No floats}: BCS excludes floating-point types, avoiding cross-platform rounding nondeterminism.
  \item \emph{Cross-language parity}: Implementations in Rust, Python, Go, JavaScript, and Move are designed to produce identical byte streams for the same schema and value.
\end{itemize}

Thus BCS induces a canonical serialization map
\[
  \sigma : \mathcal{O} \to \{0,1\}^*,
\]
from the set of operator values \( \mathcal{O} \) to byte strings.

\subsection{Operator Schemas and Hashes}

Each update law in the fractal framework is represented as a typed BCS struct. A simple Gating operator schema might be:

\begin{verbatim}
GatingOperator {
  schema_ref:   [u8; 32],  // hash of the schema record
  operator_type: u8,       // enum index for the update law class
  scale_factor: u64,       // fixed-point, no floats
  activation_fn: [u8; 32], // hash of activation descriptor
  nesting_depth: u32,      // max depth this operator applies to
}
\end{verbatim}

The canonical BCS encoding \( \sigma(o) \) of an operator instance \( o \) is then hashed to yield an operator identity hash:
\[
  H : \mathcal{O} \to \{0,1\}^{256}, \quad H(o) = \mathrm{SHA256}(\sigma(o)).
\]

\subsection{Failure of Monotonic Registry \( R \)}

A natural but flawed idea is to maintain a global, monotonic registry
\[
  R: \mathrm{Image}(H) \hookrightarrow \mathbb{P},
\]
mapping operator hashes to primes in increasing order. This registry would satisfy:

\begin{itemize}[leftmargin=*]
  \item If \( H(o) \notin \mathrm{dom}(R) \), assign the next unused prime \( p_k \) and map \( H(o) \mapsto p_k \).
  \item If \( H(o) \in \mathrm{dom}(R) \), return the existing assignment.
\end{itemize}

However, in a distributed system:

\begin{itemize}[leftmargin=*]
  \item Concurrent registrations on partitioned nodes cannot agree on the next prime without a total-order consensus protocol.
  \item Treating GitHub or a central repository as the ground truth reintroduces centralization and mutable infrastructure dependence.
\end{itemize}

Thus a monotonic \( R \) is either inconsistent or dependent on a heavy consensus layer, contradicting the stateless, trustless design goals.

\section{Topological Signatures for Fractal Trees}

\subsection{Commutative Product vs. Ordered Tree Hash}

The initial multiplicity signature for a fractal subspace \( F \) was defined as
\[
  \mathcal{M}(F) = \prod_{i=1}^k p_{n_i}^{m_i},
\]
where \( p_{n_i} \) is the prime associated to the \(i\)-th distinct operator and \( m_i \) is its multiplicity in \( F \).

This is a commutative, associative product, which means:

\begin{itemize}[leftmargin=*]
  \item Any permutation of operators yields the same signature.
  \item Nesting order and tree shape do not affect \( \mathcal{M}(F) \).
\end{itemize}

This violates the requirement that the attestation be path-dependent: different application orders of the same operators must yield distinct signatures.

\subsection{BCS-Encoded Topological Signature}

To encode the geometry and ordering of a fractal tree, we define an ordered-node BCS schema:

\begin{verbatim}
TopologicalSignature {
  operator_hash: [u8; 32],    // H(op) at this node
  children:      Vec<[u8; 32]> // ordered child signatures
}
\end{verbatim}

For a node with operator hash \( h_{op} \) and ordered child signatures \( s_1, \dots, s_k \), we encode it via BCS and define its signature:
\[
  S = \mathrm{SHA256}(\mathrm{BCS}(\texttt{TopologicalSignature})).
\]

Recursively:
\begin{itemize}[leftmargin=*]
  \item Leaf nodes: \( S_{\text{leaf}} = \mathrm{SHA256}(\mathrm{BCS}(h_{op}, [])) \).
  \item Internal nodes: \( S_{\text{node}} = \mathrm{SHA256}(\mathrm{BCS}(h_{op}, [S_{C_1},\dots,S_{C_k}])) \).
\end{itemize}

Because BCS is canonical and \( \mathrm{children} \) is an ordered sequence, swapping children changes the hash. Identical trees yield identical signatures; non-isomorphic trees (with respect to ordered children) almost surely yield different signatures.

\subsection{Illustrative Python Snippet}

Below is a simplified Python implementation of the topological signature using a hand-rolled BCS encoding for the node:

\begin{lstlisting}
import hashlib

def uleb128_encode(n: int) -> bytes:
    out = b''
    while True:
        byte = n & 0x7f
        n >>= 7
        if n:
            out += bytes([byte | 0x80])
        else:
            out += bytes([byte])
            break
    return out

def bcs_encode_topological_signature(op_hash: bytes, child_sigs: list[bytes]) -> bytes:
    if len(op_hash) != 32:
        raise ValueError("operator_hash must be 32 bytes")
    for cs in child_sigs:
        if len(cs) != 32:
            raise ValueError("child signature must be 32 bytes")
    buf = op_hash                 # operator_hash
    buf += uleb128_encode(len(child_sigs))  # children count
    for cs in child_sigs:
        buf += cs
    return buf

def topological_signature(op_hash: bytes, child_sigs: list[bytes]) -> bytes:
    encoded = bcs_encode_topological_signature(op_hash, child_sigs)
    return hashlib.sha256(encoded).digest()
\end{lstlisting}

Two trees with identical operator hashes but different child ordering produce distinct root signatures under this definition.

\section{Stateless Hash-to-Prime (H2P) Derivation}

\subsection{Motivation}

To eliminate the monotonic registry \( R \), we seek a purely functional mapping
\[
  f : \{0,1\}^{256} \to \mathbb{P},
\]
such that for each operator hash \( H_{op} \in \{0,1\}^{256} \), the assigned prime \( p_{op} = f(H_{op}) \) is:

\begin{itemize}[leftmargin=*]
  \item Deterministic and stateless.
  \item Unique with overwhelming probability: distinct operators almost surely map to distinct primes.
\end{itemize}

A naive choice is:
\[
  p_{op} = \text{NextPrime}(H_{op} \bmod 2^{64}),
\]
which uses a 64-bit truncation as a seed. This is computationally convenient but cryptographically flawed.

\subsection{Truncation Collisions and Multiplicity Collapse}

With 256-bit hashes and 64-bit truncation, many distinct hashes collide on the same seed, by the pigeonhole principle. In particular, there are \(2^{256}\) possible SHA-256 outputs but only \(2^{64}\) seeds:

\[
  \frac{2^{256}}{2^{64}} = 2^{192}
\]

On average, \(2^{192}\) operator hashes share each seed. If prime assignment is defined solely by the 64-bit seed, these distinct operators will share the same prime. This violates the multiplicity requirement:
\[
  \forall i \neq j: p_{n_i} \neq p_{n_j},
\]
and destroys the integrity of the prime-indexed multiplicity space.

\subsection{Full-Width Seeding and Deterministic Prime Generation}

To preserve the collision properties of SHA-256 and avoid truncation, we must instead use the full 256-bit hash as the seed. Conceptually:
\[
  p_{op} = G(H_{op}),
\]
where \( G \) is a deterministic prime generation algorithm mapping 256-bit seeds to 256-bit primes.

One option is to use a provable-prime generator such as Maurer's algorithm, seeded by \( H_{op} \). Maurer-type algorithms recursively construct primes of the form \( 2Rq + 1 \), using deterministic randomness and primality tests (e.g., Miller-Rabin). This yields provable primes and a seed-dependent collision profile, but introduces two problems:
\begin{itemize}[leftmargin=*]
  \item Variable and seed-dependent latency: a maliciously chosen seed can induce worst-case running time.
  \item Structural bias in the prime distribution: primes are not sampled uniformly; this affects collision estimates.
\end{itemize}

\section{Incremental Search, Offset Pinning, and First-Prime Enforcement}

\subsection{H2P with Offset Pinning}

To stabilize latency and reduce dependence on complex prime generators, we adopt a deterministic incremental search:

\begin{enumerate}[leftmargin=*]
  \item Derive an odd 256-bit integer \( N \) from \( H_{op} \) (e.g., interpret as big-endian integer and set the least significant bit to 1).
  \item Search the odd sequence \( N, N+2, N+4, \dots \) until a prime is found:
  \[
    p_{op} = N + k,
  \]
  where \( k \in 2\mathbb{Z} \) is the smallest even offset such that \( N + k \) is prime under a strong primality test (e.g., Baillie-PSW).
  \item The prover records the offset \( k \) in the execution manifest.
  \item The verifier recomputes \( N \) from \( H_{op} \), adds \( k \), and tests \( N + k \) for primality.
\end{enumerate}

This removes global state: each node can compute \( p_{op} \) locally from \( H_{op} \) and \( k \). However, two challenges remain:
\begin{itemize}[leftmargin=*]
  \item Without further constraints, a malicious prover can skip earlier primes (multiplicity forking).
  \item An unbounded search allows arbitrarily large latency.
\end{itemize}

\subsection{Offset Bound and Cramér-Type Considerations}

To bound the search, we introduce a maximum offset \( k_{\max} \), e.g.,
\[
  k_{\max} = 2^{16} = 65{,}536.
\]

Heuristically, the Cramér conjecture and related results suggest that prime gaps near \( N \approx 2^{256} \) are on the order of \( O((\log N)^2) \), which for 256 bits is roughly \( \approx 31{,}000 \). Setting \( k_{\max} \) at \( 65{,}536 \) leaves ample safety margin while guaranteeing that the prover's search terminates quickly for honest seeds.

The protocol then rejects any manifest with \( k > k_{\max} \), interpreting it as malformed or adversarial.

\subsection{First-Prime Rule and Gap Validation}

Despite offset bounding, a prover could still cheat by skipping earlier primes:

\begin{itemize}[leftmargin=*]
  \item Suppose the first prime in \( [N, N + k_{\max}] \) is at \( N + k_1 \).
  \item The prover could claim \( N + k_2 \) with \( k_2 > k_1 \) but \( k_2 \le k_{\max} \), as long as \( N + k_2 \) is also prime.
\end{itemize}

Without additional checks, the verifier would accept any such \( N+k_2 \). This breaks injectivity: the same operator hash could be associated with multiple valid primes, depending on prover behavior.

To restore determinism, the mapping must enforce a \emph{first-prime rule}:
\[
  p_{op} = \min\{ M \ge N : M \text{ is prime} \}.
\]

Enforcing this rule requires proving that every odd \( M \in [N, N+k-2] \) is composite. Two approaches were considered:

\begin{description}[leftmargin=*]
  \item[Option A (Verifier Gap Audit)] The verifier runs a primality test on every odd integer in the gap. Complexity: \( O(k) \) primality tests, which violates the \( O(1) \) verification target.
  \item[Option B (Explicit Witness Array)] The prover supplies a Fermat or Miller-Rabin witness for each odd composite in the gap; the verifier checks the witnesses. Complexity: \( O(k) \) small exponentiations and \( O(k) \) bandwidth (up to megabytes of witnesses).
\end{description}

Both options cause linear blow-up, either in verifier work or manifest size.

\section{Succinct Gap Attestation}

\subsection{Gap Circuit and zk-Gap Attestation}

To reconcile first-prime enforcement with the constant-size and \( O(1) \) verification goals, we propose a succinct proof of the gap property.

Define an arithmetic circuit \( C \) parameterized by \( N \) and \( k \) that:
\begin{itemize}[leftmargin=*]
  \item Iterates a deterministic compositeness test (e.g., Miller-Rabin with fixed bases) over all odd integers in \( [N, N+k-2] \).
  \item Outputs true if and only if each tested integer is composite.
\end{itemize}

The prover evaluates \( C \), finds the first prime at \( N+k \), and generates a succinct zero-knowledge proof \( \pi \) attesting that:

\begin{itemize}[leftmargin=*]
  \item \( C(N,k) = \text{true} \) (no prior primes in the gap).
  \item \( N+k \) itself is prime (or leaves this to an external Baillie-PSW test).
\end{itemize}

The execution manifest then includes only:
\begin{itemize}[leftmargin=*]
  \item The operator hash \( H_{op} \).
  \item The offset \( k \).
  \item The succinct proof \( \pi \).
\end{itemize}

The verifier:
\begin{enumerate}[leftmargin=*]
  \item Recomputes \( N \) from \( H_{op} \).
  \item Computes \( N+k \) and runs a single Baillie-PSW primality test.
  \item Verifies \( \pi \) via a constant-time proof system (e.g., PLONK, Halo2), typically a small number of pairing checks.
\end{enumerate}

Thus verification remains \( O(1) \) in \( k \), while the prover absorbs the linear cost of the gap computation.

\subsection{Native Multiplicity Crypto: Pedersen-Based Gap Attestation}

PhaseMirror-HQ already ships a cryptographic package, \texttt{multiplicity.crypto}, that includes:

\begin{itemize}[leftmargin=*]
  \item Production BN254 Pedersen commitments via WASM and \texttt{ark-bn254}.
  \item HMAC-SHA256-based deterministic randomness.
  \item An experimental Quantum Prime Abacus (QPA) module providing a prototype prime-coupling model and deterministic classical shadow projection.
\end{itemize}

This suggests a lighter-weight, fully native approach:

\begin{enumerate}[leftmargin=*]
  \item The prover constructs a vector of gap witnesses \( W = (w_1, \dots, w_t) \) for the composite integers in \( [N, N+k-2] \) (e.g., Fermat bases or small factorization hints).
  \item The prover computes a Pedersen commitment \(\mathsf{Com}(W)\) using BN254 via \texttt{multiplicity.crypto}'s \texttt{computeCommitment}.
  \item The manifest includes:
  \begin{itemize}
    \item \( H_{op} \), \( k \), and \( p_{op} = N+k \).
    \item The commitment \( \mathsf{Com}(W) \).
    \item Optionally, a small proof-of-knowledge or membership proof tying \( W \) to the claimed gap.
  \end{itemize}
  \item The verifier:
  \begin{itemize}
    \item Recomputes \( N \), checks the Baillie-PSW primality of \( N+k \).
    \item Verifies the commitment relation using known Pedersen verification procedures.
  \end{itemize}
\end{enumerate}

If upgraded to full zk-SNARKs over BN254 or BLS12-381, the commitment can be coupled with a succinct proof showing that \( W \) indeed corresponds to correct compositeness witnesses for all values in the gap without revealing the witnesses themselves.

\section{Formal Mathematical Overview}

We summarize the final architecture in formal terms.

\subsection{Canonical Maps and Spaces}

\begin{itemize}[leftmargin=*]
  \item Let \( \mathcal{O} \) be the set of operator values, with BCS serialization map \( \sigma : \mathcal{O} \to \{0,1\}^* \).
  \item Let \( H : \mathcal{O} \to \{0,1\}^{256} \) be \( H(o) = \mathrm{SHA256}(\sigma(o)) \).
  \item Let \( \mathcal{T} \) be the set of finite ordered rooted trees whose nodes are labeled by \( \mathcal{O} \).
\end{itemize}

\subsection{Topological Signature}

Define \( S : \mathcal{T} \to \{0,1\}^{256} \) recursively:

\begin{itemize}[leftmargin=*]
  \item For a leaf node with operator \( o \): \( S(\text{leaf}) = \mathrm{SHA256}(\mathrm{BCS}(H(o), [])) \).
  \item For a node with operator \( o \) and ordered children \( T_1,\dots,T_k \):
  \[
    S(\text{node}) = \mathrm{SHA256}(\mathrm{BCS}(H(o), [S(T_1), \dots, S(T_k)])).
  \]
\end{itemize}

\subsection{Hash-to-Prime Map}

For each operator \( o \in \mathcal{O} \):

\begin{enumerate}[leftmargin=*]
  \item Compute \( h = H(o) \).
  \item Derive seed \( N = g(h) \in \mathbb{Z}_{\ge 0} \), where \( g \) turns the 256-bit hash into an odd integer (e.g., big-endian interpretation with LSB forced to 1).
  \item For even offsets \( k \in [0, k_{\max}] \), define candidate \( M_k = N + k \).
  \item Define
  \[
    p_{op} = M_{k^\star}, \quad
    k^\star = \min \{ k \in 2\mathbb{Z}_{\ge 0} \mid M_k \text{ is prime} \}.
  \]
\end{enumerate}

If no prime is found with \( k \le k_{\max} \), the operator is invalid or the system parameters must be adjusted.

\subsection{Gap Attestation Relation}

For given \( (h, k, p) \), define the gap property:
\[
  \mathsf{Gap}(h,k,p) \equiv
  \left(
    p = g(h) + k
  \right)
  \land
  \left(
    \forall \ell \in 2\mathbb{Z}_{\ge 0}, 0 \le \ell < k :
    g(h)+\ell \text{ is composite}
  \right)
  \land
  \left(
    p \text{ is prime}
  \right).
\]

The prover generates a proof \( \pi \) that \( \mathsf{Gap}(h,k,p) \) holds, using either:
\begin{itemize}[leftmargin=*]
  \item A zk-SNARK over an arithmetic circuit encoding the primality/compositeness checks; or
  \item A Pedersen commitment and auxiliary relations under \texttt{multiplicity.crypto}.
\end{itemize}

The verifier accepts the operator-prime mapping if and only if:

\begin{itemize}[leftmargin=*]
  \item The cryptographic gap proof \( \pi \) verifies.
  \item \( p \) passes a strong primality test (e.g., Baillie-PSW) at verification time.
\end{itemize}

Given the collision-resistance of SHA-256 and the statistical distribution of primes, the probability that two distinct operators yield the same prime is negligible (far below \( 2^{-128} \)), and the mapping \( o \mapsto p_{op} \) is injective with overwhelming probability.

\section{Integration with \texttt{multiplicity.crypto}}

The PhaseMirror-HQ repository includes a \texttt{multiplicity.crypto} package with:

\begin{itemize}[leftmargin=*]
  \item A production BN254 Pedersen commitment implementation via WASM and Rust (\texttt{commitment.rs}, \texttt{commitment.ts}).
  \item HMAC-SHA256-based deterministic randomness in \texttt{randomness.ts}.
  \item A Python bridge in \texttt{multiplicity/crypto/python}, exporting functions like \texttt{computeCommitment}, \texttt{merkleRoot}, and \texttt{getBridgeStatus}.
  \item An experimental QPA module in \texttt{qpa.ts} for modeling prime coupling.
\end{itemize}

The recommended path to production is:

\begin{enumerate}[leftmargin=*]
  \item Promote the QPA prime-coupling model (\texttt{multiplicity.crypto.qpa}) to a stable H2P interface: \texttt{derivePrimeFromHash(H\_op) -> p\_op, k}.
  \item Implement the gap-attestation commitment using the existing Pedersen commitment API: \texttt{computeCommitment(message, randomness)}.
  \item Use \texttt{merkleRoot} to build the fractal Merkle structure from the topological signatures.
  \item Provide Python and TypeScript bindings that wrap the entire attestation pipeline into a single call for the higher-level framework.
\end{enumerate}

\section{Code Appendix: Reference H2P Implementation}

\subsection{Python H2P Skeleton}

\begin{lstlisting}
import hashlib

def hash_operator_to_seed(op_bcs: bytes) -> int:
    """Full-width seed N from BCS-encoded operator."""
    h = hashlib.sha256(op_bcs).digest()
    N = int.from_bytes(h, "big")
    # Force N to be odd
    if N % 2 == 0:
        N += 1
    return N

def is_probable_prime(n: int) -> bool:
    """Placeholder: Baillie-PSW or equivalent."""
    # Implementation omitted; use a proven BPSW library in production
    ...

def next_prime_with_offset(N: int, k_max: int = 2**16) -> tuple[int, int]:
    """Deterministic incremental search with offset bounding."""
    for k in range(0, k_max+1, 2):  # step by 2 over odd candidates
        candidate = N + k
        if is_probable_prime(candidate):
            return candidate, k
    raise ValueError("No prime found within k_max")

def derive_prime_from_operator(op_bcs: bytes, k_max: int = 2**16) -> tuple[int, int]:
    N = hash_operator_to_seed(op_bcs)
    p, k = next_prime_with_offset(N, k_max)
    return p, k
\end{lstlisting}

In a production implementation, \texttt{is\_probable\_prime} would use a high-quality Baillie-PSW library. The gap attestation proof (either zk-SNARK or Pedersen-based) is layered on top of this core.

\section{Conclusion and Future Work}

This report presents a complete, stateless architecture for mapping operator hashes to unique primes in a prime-indexed multiplicity space. The design stands on four pillars:

\begin{enumerate}[leftmargin=*]
  \item BCS-based canonical serialization ensures cross-language determinism.
  \item Topological signatures provide path-dependent, non-commutative invariants for fractal trees.
  \item H2P with offset pinning and first-prime enforcement gives each operator a unique prime without global state.
  \item Native \texttt{multiplicity.crypto} primitives (Pedersen commitments, HMAC-SHA256 randomness, QPA) anchor the design in a production-ready cryptographic substrate.
\end{enumerate}

Future work includes:

\begin{itemize}[leftmargin=*]
  \item Formalizing the QPA prime-coupling model as a certified H2P library.
  \item Benchmarking Baillie-PSW and gap-attestation proof generation at scale.
  \item Extending the BCS schema registry to support schema evolution with explicit cryptographic lineage.
  \item Integrating BLS12-381-based zk-SNARKs or Halo2-based proofs for fully succinct, zero-knowledge gap attestation.
\end{itemize}

Taken together, these developments provide a rigorous, cryptographically grounded foundation for prime-indexed multiplicity in the PhaseMirror fractal framework.

\appendix
\section{Mathematical Appendix}
\label{app:math}

This appendix collects the formal statements and proofs that underlie the constructions
described in Sections~\ref{sec:canonical}--\ref{sec:integration}.  
All symbols are defined in the main paper; we repeat the most important ones for
convenience.

\subsection{Notation}
\begin{itemize}
  \item $\mathcal{O}$ – the set of all operator values (instances of the BCS‑encoded
        structs that describe update laws).
  \item $\sigma:\mathcal{O}\rightarrow\{0,1\}^{*}$ – the Binary Canonical Serialization
        map (fixed field order, ULEB128‑prefixed sequences, no floats) \cite{web:46}.
  \item $H:\mathcal{O}\rightarrow\{0,1\}^{256}$ – the operator hash,
        $H(o)=\operatorname{SHA256}(\sigma(o))$.
  \item $g:\{0,1\}^{256}\rightarrow\mathbb{Z}_{\ge0}$ – the seed‑translation function
        (big‑endian integer interpretation of the hash, with the least significant bit
        forced to 1 so that $g(h)$ is odd).
  \item $k_{\max}=2^{16}=65536$ – the maximal allowed offset (see
        Section~\ref{subsec:offset}).
  \item $\mathbb{P}$ – the ordered set of primes $\{2,3,5,7,\dots\}$.
  \item For $h\in\{0,1\}^{256}$ we write $N=g(h)$ and
        $p(h,k)=N+k$ (the candidate prime at offset $k$).
\end{itemize}

\subsection{Lemma 1 (Seed Range)}
For every operator hash $h\in\{0,1\}^{256}$ the seed $N=g(h)$ satisfies
\[
2^{255}\le N<2^{256}.
\]
\begin{proof}
  By construction $g$ interprets the 256‑bit string as an unsigned integer in
  $[0,2^{256})$.  Forcing the least significant bit to 1 can only increase the value,
  therefore $N\ge 2^{255}$ (the smallest odd number with the most‑significant bit set).
  The upper bound follows directly from the 256‑bit width.  ∎
\end{proof}

\subsection{Lemma 2 (Existence of a Prime within the Offset Bound)}
Let $N$ be an odd integer with $2^{255}\le N<2^{256}$.
Then there exists an even $k$ with $0\le k\le k_{\max}$ such that $p=N+k$ is prime.
\begin{proof}
  The maximal known prime gap for numbers below $2^{66}$ is $1724$ \cite{web:92}.
  Heuristic and conditional results (Cramér’s conjecture) predict that for numbers
  of size $x$ the maximal gap $G(x)$ satisfies $G(x)=O((\log x)^2)$.
  With $x\approx2^{256}$ we have $\log x\approx256\ln2\approx177$, hence
  $(\log x)^2\approx3.1\times10^{4}<k_{\max}=65536$.
  Thus any interval of length $k_{\max}$ that starts at an odd $N\ge2^{255}$ must
  contain at least one prime.  The statement holds unconditionally for all
  numbers $<2^{66}$ by the explicit table of maximal gaps \cite{web:91},
  and for larger numbers it follows from the widely accepted Cramér‑type bound,
  which is assumed in the design of the offset limit.  ∎
\end{proof}

\subsection{Lemma 3 (Collision Probability of Full‑Width Seeding)}
Let $h_1\neq h_2\in\{0,1\}^{256}$ be two distinct operator hashes and
let $p_i=p(g(h_i),k_i)$ be the primes obtained by the deterministic incremental
search with offset bound $k_{\max}$.  Then
\[
\Pr\bigl[p_1=p_2\bigr]\le 2^{-248}.
\]
\begin{proof}
  The map $h\mapsto N=g(h)$ is injective (adding the fixed LSB does not cause
  collisions).  Hence distinct hashes give distinct odd seeds $N_1\neq N_2$.
  The search algorithm returns the \emph{first} prime at or after the seed.
  Therefore $p_i$ is the smallest prime $\ge N_i$.
  If $p_1=p_2=p$, then both seeds lie in the interval
  $[p-k_{\max},\,p]$ and are distinct odd numbers.
  The interval length is $k_{\max}=2^{16}$, so it contains at most
  $2^{15}$ odd numbers.
  By the prime number theorem the density of primes near $2^{256}$ is
  $\approx 1/\ln(2^{256})\approx1/177$.
  Hence the expected number of primes in the interval is at most
  $\displaystyle\frac{2^{15}}{177}<2^{8}=256$.
  Consequently the probability that two independently chosen seeds fall into the
  same prime‑generating interval is bounded by
  $\frac{256}{2^{256}}=2^{-248}$.
  (A more precise bound uses the exact count of primes in $[2^{255},2^{256})$,
  which is $>2^{248}$; the inequality above is therefore conservative.)  ∎
\end{proof}

\subsection{Theorem 1 (Injectivity of the Hash‑to‑Prime Map)}
Define the map
\[
\Phi:\mathcal{O}\longrightarrow\mathbb{P},\qquad
\Phi(o)=p\bigl(g(H(o)),k^\star(o)\bigr)
\]
where $k^\star(o)$ is the smallest even offset such that
$p\bigl(g(H(o)),k^\star(o)\bigr)$ is prime (or $\bot$ if no such offset exists
within $k_{\max}$).  Then, with probability at least $1-2^{-248}$ over the
choice of the operator hash $H(o)$, $\Phi$ is injective on $\mathcal{O}$.
\begin{proof}
  Assume two distinct operators $o_1\neq o_2$ satisfy $\Phi(o_1)=\Phi(o_2)=p\neq\bot$.
  Let $h_i=H(o_i)$ and $N_i=g(h_i)$.  By definition of $k^\star$ we have
  $p=N_i+k_i$ with $k_i$ the smallest even offset yielding a prime.
  Since $p$ is the same prime, both seeds lie in the interval
  $[p-k_{\max},\,p]$ and are distinct odd numbers (otherwise $h_1=h_2$).
  The argument of Lemma~3 then applies, giving
  $\Pr[\text{collision}]\le2^{-248}$.
  Hence, except with that negligible probability, distinct operators map to
  distinct primes.  ∎
\end{proof}

\subsection{Corollary 1 (Operator‑Norm Bound)}
If we view the mapping $\Phi$ as an operator from the Hilbert space
$\ell^{2}(\mathcal{O})$ (with the counting measure on $\mathcal{O}$) to
$\ell^{2}(\mathbb{P})$ defined by
\[
(\Phi f)(p)=\sum_{o\in\Phi^{-1}\{p\}}f(o),
\]
then $\|\Phi\|_{2\to2}\le\sqrt{1+2^{-248}}<1.0000001$.
\begin{proof}
  For any $f\in\ell^{2}(\mathcal{O})$,
  \[
    \|\Phi f\|_{2}^{2}
    =\sum_{p\in\mathbb{P}}\Bigl|\sum_{o\in\Phi^{-1}\{p\}}f(o)\Bigr|^{2}
    \le\sum_{p\in\mathbb{P}}\sum_{o\in\Phi^{-1}\{p\}}|f(o)|^{2}
      \;+\!\!\sum_{p}\!\!\sum_{\substack{o_1\neq o_2\\ \Phi(o_1)=\Phi(o_2)=p}}
            |f(o_1)||f(o_2)|.
  \]
  The first term equals $\|f\|_{2}^{2}$ because the preimages $\Phi^{-1}\{p\}$ are
  disjoint.  The second term counts ordered pairs of distinct operators that
  collide under $\Phi$; by Theorem~1 the expected number of such pairs is at most
  $2^{-248}\|f\|_{2}^{2}$ (each collision contributes at most $|f(o_1)||f(o_2)|
  \le\frac12(|f(o_1)|^{2}+|f(o_2)|^{2})$).  Hence
  $\|\Phi f\|_{2}^{2}\le(1+2^{-248})\|f\|_{2}^{2}$, giving the claimed norm bound. ∎
\end{proof}

\subsection{Remark on Operator Norm}
The bound above shows that $\Phi$ is essentially an isometry on $\ell^{2}$:
it does not amplify energy by more than a factor of $1+2^{-248}$, which is
far below any realistic precision threshold.  Consequently, the hash‑to‑prime
step can be treated as norm‑preserving in practical analyses of the
fractional‑update dynamics of the PhaseMirror system.

\section*{References (used in the appendix)}
\begin{thebibliography}{99}
  \bibitem{web:46}
    aptos-labs/bcs, ``Binary Canonical Serialization Library'', GitHub,
    \url{https://github.com/aptos-labs/bcs}, accessed 2026-04-24.

  \bibitem{web:37}
    ``Why should hash functions use a prime number modulus?'',
    StackOverflow, \url{https://stackoverflow.com/questions/1145217/why-should-hash-functions-use-a-prime-number-modulus},
    accessed 2026-04-24.

  \bibitem{web:91}
    Table of Known Maximal Gaps, The Prime Pages,
    \url{http://primerecords.dk/primegaps/megagap3.htm}, accessed 2026-04-24.

  \bibitem{web:92}
    Prime gap – Wikipedia,
    \url{https://en.wikipedia.org/wiki/Prime_gap}, accessed 2026-04-24.

  \bibitem{cite:155}
    multiplicity.crypto README, PhaseMirror‑HQ,
    \url{https://raw.githubusercontent.com/MultiplicityFoundation/PhaseMirror-HQ/main/multiplicity/src/multiplicity/crypto/README.md},
    accessed 2026-04-24.

  \bibitem{cite:156}
    multiplicity.crypto source directory, PhaseMirror‑HQ,
    \url{https://github.com/MultiplicityFoundation/PhaseMirror-HQ/tree/main/multiplicity/src/multiplicity/crypto},
    accessed 2026-04-24.
\end{thebibliography}


\nocite{*}
\bibliographystyle{unsrt}  
\bibliography{references}
\end{document}